% Preámbulo común del documento de diseño y del informe técnico de PEDP
\usepackage{instructivo}
\renewcommand{\IEEEkeywordsname}{Palabras clave}
\usepackage[american,siunitx]{circuitikz}

\sisetup{per-mode=symbol, group-digits=integer, output-decimal-marker={,}}
\graphicspath{{./}{./fig/}}

\usepackage{amsmath}
\usepackage{float}
\usepackage[natbib=true,style=trad-unsrt,backend=biber]{biblatex}
\addbibresource{literatura.bib}
\renewcommand*{\bibfont}{\footnotesize}

% Columnas de tabularx que parten el texto en varias líneas para no salirse de la columna
\newcolumntype{L}{>{\raggedright\arraybackslash}X}
\newcolumntype{C}{>{\centering\arraybackslash}X}

% Marcador visible para datos que todavía no existen, nunca se llena con un valor calculado
\newcommand{\pendiente}[1]{\textcolor{red}{\textbf{[#1]}}}

\hyphenation{op-tical net-works semi-conduc-tor}
